\noindent
{\Large \textbf{\section{Математическое описание}}} \par

Кэш-память хранит копии данных основной памяти и делится на уровни (L1, L2, L3), каждый следующий из которых больше и медленнее предыдущего. Данные переносятся из памяти в кэш не по одному байту, а блоками фиксированного размера $L$ - строками кэша. Время обращения к данным зависит от того, в каком уровне они находятся, поэтому по времени можно судить о характеристиках кэша, не обращаясь к документации процессора.

\vspace{6pt}
\textbf{Среднее время обращения.} В эксперименте измеряется полное время $T$ выполнения $N$ обращений к памяти, а среднее время одного обращения вычисляется как
\[
    t = \frac{T}{N}.
\]
Деление на большое $N$ необходимо, так как одно обращение занимает доли наносекунды и не может быть измерено системным таймером непосредственно.

\vspace{6pt}
\textbf{Математическое ожидание результатов замеров.} На время одного замера влияют различные помехи (из-за планировщика операционной системы), поэтому замер повторяется $n$ раз, и результатом считается математическое ожидание полученных значений $t_1, \ldots, t_n$. Все замеры считаются равновероятными, то есть вероятность каждого исхода равна $1/n$, и тогда
\[
    E[t] = \sum_{j=1}^{n} \frac{1}{n} \, t_j = \frac{1}{n} \sum_{j=1}^{n} t_j.
\]
В работе используется $n = 10$ измерений для каждой величины.

\vspace{8pt}
{\Large \textbf{\subsection{Размер кэша}}} \par
Пусть рабочий массив имеет размер $S$, а размер уровня кэша равен $C$. Если все элементы массива многократно обходятся, то при $S \leqslant C$ все данные находятся в этом уровне и $t$ равно задержке этого уровня. При $S > C$ часть данных вытесняется и берется из следующего, более медленного уровня, поэтому $t$ возрастает. Размером уровня считается наибольший $S$, при котором время еще соответствует этому уровню, за ним следует скачок. Для трех уровней в зависимости $t(S)$ ожидаются три плато, разделенные скачками.

\vspace{8pt}
{\Large \textbf{\subsection{Размер строки}}} \par
Пусть массив обходится с шагом $d$ байт и значительно превышает размер кэша, так что строка загружается из основной памяти. Пока $d < L$, в одну загруженную строку попадают $L / d$ последовательных обращений, из которых только первое требует загрузки. При $d \geqslant L$ каждое обращение попадает в новую строку и требует загрузки. Следовательно, доля обращений с загрузкой равна $\min(1, d / L)$, и время $t(d)$ возрастает с ростом $d$ до значения $d = L$, после чего перестает расти по этой причине. Размером строки считается шаг, на котором наблюдается наибольшее увеличение $t$ относительно предыдущего шага.

\vspace{8pt}
{\Large \textbf{\subsection{Ассоциативность}}} \par
Множественно-ассоциативный кэш ($k$-way set associative) разделен на наборы (sets), в каждом из которых $k$ мест (путей, ways). Блок памяти может находиться только в одном определенном наборе, номер которого определяется адресом, но внутри набора занимает любое свободное место. При размере кэша $C$ и строке $L$ число наборов равно
\[
    m = \frac{C}{k \cdot L}.
\]
Два адреса, различающиеся на величину, кратную $C / k$, попадают в один набор. Если взять $n$ ячеек, расположенных с таким шагом, то все они борются за $k$ мест одного набора. При $n \leqslant k$ все ячейки одновременно находятся в кэше и время $t$ мало, а при $n = k + 1$ хотя бы одна ячейка вытесняется, и $t$ возрастает. Следовательно, ассоциативность равна наибольшему $n$, при котором время еще не увеличилось.

\newpage
